\section{Flujo inicial de navegaci\'on}
% TODO: Describir el recorrido principal del usuario por las pantallas.

\subsection{Diagrama de navegaci\'on}
% TODO: Reemplazar el diagrama de ejemplo por el flujo real del sistema.
\begin{figure}[H]
  \centering
  \begin{tikzpicture}[
      nodo/.style={rectangle, draw=black, rounded corners, minimum width=3cm,
                   minimum height=1cm, align=center, fill=gray!12},
      flecha/.style={-{Stealth[length=3mm]}, thick}
    ]
    \node[nodo] (inicio) {Inicio};
    \node[nodo, right=1.5cm of inicio] (mapa) {Mapa interactivo};
    \node[nodo, below=1.5cm of mapa] (detalle) {Detalle del recinto};
    \draw[flecha] (inicio) -- (mapa);
    \draw[flecha] (mapa) -- (detalle);
  \end{tikzpicture}
  \caption{Flujo de navegaci\'on propuesto (ejemplo)}
\end{figure}
